\section{General-family continuation: uniform mixing bounds and the original kernel determinant}\label{general-family-continuation-uniform-mixing-bounds-and-the-original-kernel-determinant}

\textbf{The new GitHub calculation can now be combined with a stronger, dimension-controlled mixing theorem.} For the original covariance activation

\[
G(t)=(G^{-1}+t\Omega)^{-1},
\]

the complete motion of the observation's minimum section satisfies

\[
\boxed{
\int_0^\infty t\,\|\dot L(t)\|_{\mathrm{HS}(Q(t)\to G(t))}^{2}\,dt
\le \operatorname{rank}(\Lambda\Omega\Lambda^*)
\le \dim\operatorname{im}\Lambda.
}
\tag{1}
\]

No covariance eigenvalue separation, equality of primary lengths, or bound on the condition number is required.

On the original simple-quartet observation domain, this gives

\[
\boxed{
\int_0^1t\,\|\dot L_N(t)\|_{\mathrm{HS}}^2\,dt
\le q-8k+16.
}
\tag{2}
\]

This improves the preceding \(O_h(kq\log(q/k))\) bound for that integral to an explicit \(O(q)\) bound. It concerns the \textbf{actual minimum-section motion}, not a substitute norm or an independently chosen section.

The generalization also reaches every fixed flag quotient: its determinant change has a positive spectral-shift representation, and all its derivatives are bounded by the quotient's rank. For several noncommuting source additions, the large-parameter determinant exponents are determined by an explicit filtration of their actual value ranges.

There is a necessary boundary correction. A source family can lose rank even while the supplied direction-preserving coercivity inequality remains exact. I calculate that phenomenon below, including the precise additional minimum needed to continue the family across such a point.

I incorporated the new original-kernel proof at \textbf{\nolinkurl{5992ad50c0f2a06921f1fcaded7b4e38097c0739}}. The final recheck returned \textbf{\nolinkurl{487253b5b01851da5b9b98af4e146fecf1fab902}}; its changes are publication links and metadata, not changes to the mathematical proof used here. The new source itself expressly leaves the original kernel's large-degree \(kq\) coefficient unevaluated. I do not relabel its exact determinant formula as that coefficient.

\subsection{1. Retain the original mixed source and its complete minimum}\label{retain-the-original-mixed-source-and-its-complete-minimum}

At one original cutoff, retain

\[
H_Z=\Gamma-C_X^*H_X^{-1}C_X,\qquad
F_Z=I-J_NH_X^{-1}C_X,
\]

\[
\Omega=F_ZH_Z^+F_Z^*,\qquad
G^J=(G^{-1}+\Omega)^{-1}.
\]

The inverse \(H_Z^+\) acts on its entire positive range. The original normal representative proves

\[
F_Z^*G_{k,1}F_Z\preceq \Lambda_{h,n}^{\,k}H_Z,
\]

so the value map vanishes on the kernel of the normal Gram. These are the supplied full cross-Gram identities, including their interference terms.

Let

\[
\Lambda:E\longrightarrow B=\operatorname{im}\Lambda
\]

be the original observation, with

\[
K=\ker\Lambda,\qquad m=\dim K,\qquad r=\dim B,\qquad q=m+r.
\]

All results through §5 use these \textbf{actual dimensions}. The special formula \(m=8k-16\) is used only on its proved simple-quartet five-orbit domain.

Write

\[
H_0=I_K^*GI_K,\qquad
Q_0=(\Lambda G^{-1}\Lambda^*)^{-1},\qquad
L_0=G^{-1}\Lambda^*Q_0.
\]

The explicit isometry and its inverse are

\[
F=[\,I_KH_0^{-1/2},\,L_0Q_0^{-1/2}\,],
\qquad
F^*GF=I,\qquad F^{-1}=F^*G.
\tag{3}
\]

In these coordinates,

\[
Q_0^{1/2}\Lambda F=[0,I_r].
\]

Transport the actual added covariance by that same map:

\[
\mathcal K=F^{-1}\Omega F^{-*}
=
\begin{pmatrix}
U&V\\V^*&W
\end{pmatrix}\succeq0.
\tag{4}
\]

For \(t\ge0\), put

\[
D_t=I+tW,\qquad
X_t=tVD_t^{-1},
\]

\[
E_t=I+tU-t^2VD_t^{-1}V^*.
\tag{5}
\]

Full block inversion gives

\[
F^*G(t)F=(I+t\mathcal K)^{-1},
\]

and consequently

\[
\boxed{
\begin{aligned}
I_K^*G(t)I_K&=H_0^{1/2}E_t^{-1}H_0^{1/2},\\
Q(t)&=Q_0^{1/2}D_t^{-1}Q_0^{1/2},\\
L(t)&=F\binom{X_t}{I_r}Q_0^{1/2}.
\end{aligned}}
\tag{6}
\]

Thus the derivative in (1) is exactly the derivative of the original observation minimum section. It is not the derivative of an arbitrarily transported frame.

The original arithmetic action changes by the kernel leakage already calculated in the pasted source:

\[
A_B^J-A_B
=\Lambda A P_K\Omega\Lambda^*Q^J.
\]

The analysis below retains this term throughout.

\section{2. The entire kernel determinant has a positive spectral-shift representation}\label{the-entire-kernel-determinant-has-a-positive-spectral-shift-representation}

Define its actual logarithmic loss

\[
f_K(t)
=
\log\det H_0-\log\det(I_K^*G(t)I_K).
\]

Equation (6) gives

\[
\boxed{
f_K(t)
=\log\det E_t
=\log\det(I+t\mathcal K)-\log\det(I+tW).
}
\tag{7}
\]

Let

\[
\xi_K(s)
=
\#\{\lambda(\mathcal K)>s\}
-\#\{\lambda(W)>s\},
\qquad s>0,
\]

counting multiplicities. The map \(v\mapsto(0,v)\) makes \(W\) the literal compression of \(\mathcal K\). Min--max therefore proves

\[
0\le\xi_K(s)\le m.
\tag{8}
\]

Since

\[
\log(1+t\lambda)
=\int_0^\lambda\frac{t}{1+ts}\,ds,
\]

we obtain the exact representation

\[
\boxed{
f_K(t)=
\int_0^\infty\frac{t}{1+ts}\,\xi_K(s)\,ds.
}
\tag{9}
\]

Every zero covariance eigenvalue is retained: it simply contributes an empty interval in this expression.

For every integer \(n\ge1\),

\[
\boxed{
(-1)^{n-1}f_K^{(n)}(t)
=
n!\int_0^\infty
\frac{s^{n-1}\xi_K(s)}{(1+ts)^{n+1}}\,ds,
}
\]

and hence

\[
\boxed{
0\le(-1)^{n-1}f_K^{(n)}(t)
\le \frac{m(n-1)!}{t^n},
\qquad t>0.
}
\tag{10}
\]

This simultaneously controls monotonicity, concavity, and every higher derivative. The controlling integer is the original kernel rank, not the full packet dimension.

Representation (9) is an explicit complete-Bernstein representation. That function class has an established matrix-valued theory; the result used here is proved directly from the two actual covariance matrices and their compression.

\subsection{2.1 The large-activation exponent and constant are evaluated}\label{the-large-activation-exponent-and-constant-are-evaluated}

Positivity of (4) proves

\[
\ker W\subset\ker V.
\]

Define the full shorted block

\[
S=U-VW^+V^*\succeq0.
\tag{11}
\]

The determinant-one block congruence

\[
\begin{pmatrix}I&-VW^+\\0&I\end{pmatrix}
\mathcal K
\begin{pmatrix}I&0\\-W^+V^*&I\end{pmatrix}
=
\begin{pmatrix}S&0\\0&W\end{pmatrix}
\]

gives

\[
a:=\operatorname{rank}S
=\operatorname{rank}\mathcal K-\operatorname{rank}W.
\]

Using the actual nonzero eigenvalues in (7),

\[
\boxed{
f_K(t)
=
a\log t+
\log\frac{\operatorname{pdet}\mathcal K}
{\operatorname{pdet}W}
+R(t),
}
\tag{12}
\]

with the finite bound

\[
\boxed{
|R(t)|
\le
\frac{\operatorname{Tr}\mathcal K^+
+\operatorname{Tr}W^+}{t}.
}
\tag{13}
\]

The pseudodeterminant of a zero-dimensional positive range is the empty product, one.

There is an expression for the constant that displays the mixing left when \(S\) is singular. Let \(P_0\) be the orthogonal projection onto \(\ker S\), and \(S_+\) the restriction to \(\operatorname{ran}S\). Then

\[
\boxed{
\frac{\operatorname{pdet}\mathcal K}{\operatorname{pdet}W}
=
\det S_+\,
\det_{\ker S}\!\left(
I+P_0VW^{+2}V^*P_0
\right).
}
\tag{14}
\]

To prove it, the exact identity

\[
E_t
=
I+tS+tVW^+(I+tW)^{-1}V^*
\]

gives

\[
E_t=tS+B_\infty+O(t^{-1}),
\qquad
B_\infty=I+VW^{+2}V^*.
\]

Block factorization on \(\operatorname{ran}S\oplus\ker S\) produces

\[
\det E_t
=
t^a\det S_+\det(P_0B_\infty P_0|_{\ker S})
[1+O(t^{-1})].
\]

Comparison with (12) proves (14). In particular, the extra determinant on \(\ker S\) is not omitted by treating a pseudodeterminant as multiplicative under a nonunitary congruence.

The full metric scales are also determined. With

\[
J_t=\operatorname{diag}(\sqrt t\,I_a,I_{m-a}),
\]

\[
\boxed{
J_tE_t^{-1}J_t
\longrightarrow
\operatorname{diag}
\left(S_+^{-1},
(P_0B_\infty P_0|_{\ker S})^{-1}\right).
}
\tag{15}
\]

Thus exactly \(a\) kernel metric directions have squared-norm scale \(t^{-1}\); the remaining directions approach the displayed positive limiting metric.

\subsection{2.2 The slope counts actual unobserved normal values}\label{the-slope-counts-actual-unobserved-normal-values}

Factor the original covariance through its original normal source:

\[
F^{-1}F_ZH_Z^{+1/2}
=
\binom{B_K}{B_B}.
\]

Then

\[
W=B_BB_B^*,\qquad
V=B_KB_B^*,
\]

and

\[
\boxed{
S=B_KP_{\ker B_B}B_K^*.
}
\tag{16}
\]

The exponent \(a\) in (12) therefore counts the independent kernel values obtainable from normal source directions whose observed value is zero. This is an exact value-map statement. It does not follow merely from the dimension of \(H_Z\).

\section{3. The section's total motion is bounded independently of all covariance scales}\label{the-sections-total-motion-is-bounded-independently-of-all-covariance-scales}

Differentiating (5) gives

\[
X_t'=VD_t^{-2},
\qquad
E_t''=-2VD_t^{-3}V^*.
\]

The motion energy in the original quotient norms is

\[
\boxed{
\mathcal E(t)
:=
\|\dot L(t)\|_{\mathrm{HS}(Q(t)\to G(t))}^{2}
=
\operatorname{Tr}(E_t^{-1}VD_t^{-3}V^*).
}
\tag{17}
\]

The complete curvature identity is

\[
\boxed{
-f_K''(t)
=
\operatorname{Tr}
\left[(E_t^{-1/2}E_t'E_t^{-1/2})^2\right]
+2\mathcal E(t).
}
\tag{18}
\]

This is the metric-drift term plus the actual minimum-section-motion term.

\subsection{3.1 Resolve all positive covariance modes without assuming they commute with the kernel blocks}\label{resolve-all-positive-covariance-modes-without-assuming-they-commute-with-the-kernel-blocks}

Let \(P_\lambda\) be the spectral projections of \(W\) for its distinct positive eigenvalues. Define

\[
B_\lambda=\lambda^{-2}VP_\lambda V^*\succeq0.
\]

The matrices \(B_\lambda\) need not commute.

The exact Schur expression becomes

\[
\boxed{
E_t=I+tS+
\sum_{\lambda>0}\frac{t\lambda}{1+t\lambda}B_\lambda,
}
\tag{19}
\]

and

\[
\mathcal E(t)
=
\sum_{\lambda>0}
\frac{\lambda^2}{(1+t\lambda)^3}
\operatorname{Tr}(E_t^{-1}B_\lambda).
\tag{20}
\]

For the pointwise bound, use

\[
\frac{(t\lambda)^2}{(1+t\lambda)^3}
\le\frac14\frac{t\lambda}{1+t\lambda}.
\]

Summing inside the trace gives

\[
\boxed{
t^2\mathcal E(t)
\le
\frac14\operatorname{Tr}
\left[E_t^{-1}(E_t-I-tS)\right]
\le\frac m4.
}
\tag{21}
\]

For the full integral, retain one \(\lambda\) term at a time in (19):

\[
E_t\succeq I+\frac{t\lambda}{1+t\lambda}B_\lambda.
\]

Set \(s=t\lambda/(1+t\lambda)\). The corresponding integral is bounded by

\[
\int_0^1
\operatorname{Tr}\bigl[sB_\lambda(I+sB_\lambda)^{-1}\bigr]\,ds.
\]

For an eigenvalue \(b>0\), that scalar integral is

\[
1-\frac{\log(1+b)}b.
\]

Define

\[
\varphi(b)=
\begin{cases}
1-\log(1+b)/b,&b>0,\\
0,&b=0.
\end{cases}
\]

Then

\[
\boxed{
\int_0^\infty t\mathcal E(t)\,dt
\le
\sum_{\lambda>0}\operatorname{Tr}\varphi(B_\lambda)
\le
\sum_{\lambda>0}\operatorname{rank}(VP_\lambda)
\le\operatorname{rank}W.
}
\tag{22}
\]

This proves (1). For a finite upper endpoint \(T\), each positive eigenvalue \(b\) contributes the sharper bound

\[
s_T-\frac{\log(1+s_Tb)}b,
\qquad
s_T=\frac{T\lambda}{1+T\lambda}.
\tag{23}
\]

A complementary bound retains the complete coupling determinant:

\[
\boxed{
\int_0^\infty t\mathcal E(t)\,dt
\le
\frac12\log\det(I+VW^{+2}V^*).
}
\tag{24}
\]

Indeed, dropping only the nonnegative \(tS\) enlarges the inverse in the energy. For the resulting \(E_t^{(0)}\), integrate (18); its limit is \(I+VW^{+2}V^*\), and \(t(\log\det E_t^{(0)})'\to0\).

Equations (22) and (24) may be used together by taking their minimum.

\subsection{3.2 The observed-dimension constant is sharp}\label{the-observed-dimension-constant-is-sharp}

The bound cannot generally be replaced by the kernel dimension.

Take a one-dimensional kernel and an \(r\)-dimensional observation. For \(R>1\), set

\[
W_R=\operatorname{diag}(R^{-2},R^{-4},\ldots,R^{-2r}),
\]

\[
V_{R,j}=R^{-3j/2},
\qquad
U_R=V_RW_R^{-1}V_R^*.
\]

The full covariance is positive semidefinite and \(S=0\). Its scalar \(B_{\lambda_j}\) equals \(R^j\).

At times \(t=u/\lambda_j\), the \(j\)-th summand dominates \(E_t\) as \(R\to\infty\), while the earlier modes have already saturated and the later modes remain smaller. On every fixed \(0<a<b<\infty\), the \(j\)-th energy integral over that interval tends to

\[
\int_a^b\frac{du}{(1+u)^2}.
\]

The \(r\) intervals are disjoint for large \(R\). Letting \(a\downarrow0\), \(b\uparrow\infty\), and using (22), proves

\[
\boxed{
\int_0^\infty t\mathcal E_R(t)\,dt\longrightarrow r.
}
\tag{25}
\]

This is an explicit general-family example, not an arithmetic packet substituted into the programme.

\subsection{3.3 The finite motion does not erase determinant growth}\label{the-finite-motion-does-not-erase-determinant-growth}

For a single kernel and observation mode, take

\[
W=1,\qquad V=\sqrt\rho,\qquad U=\rho,\qquad \rho>0.
\]

Then

\[
f_K(t)
=
\log\frac{1+(1+\rho)t}{1+t},
\]

\[
\mathcal E(t)
=
\frac{\rho}{(1+t)^2[1+(1+\rho)t]}.
\]

The two quantities evaluate to

\[
\boxed{
\int_0^\infty t\mathcal E(t)\,dt
=
1-\frac{\log(1+\rho)}{\rho}<1,
}
\]

\[
\boxed{f_K(\infty)=\log(1+\rho).}
\tag{26}
\]

The determinant loss can therefore be arbitrarily large while the complete weighted section motion stays below one. Equation (18) locates the remaining change in the metric-drift term.

More generally, combining (12), (18), and (22) gives

\[
\int_0^Tt\,
\operatorname{Tr}
[(E_t^{-1/2}E_t'E_t^{-1/2})^2]\,dt
=
a\log T+O(1).
\tag{27}
\]

The coefficient \(a\) is the actual rank in (16).

\subsection{3.4 Return the estimate through the arithmetic action}\label{return-the-estimate-through-the-arithmetic-action}

Write

\[
F^{-1}AF=
\begin{pmatrix}
A_{KK}&A_{KB}\\
A_{BK}&A_{BB}
\end{pmatrix}.
\]

The observed action in the fixed value coordinates is

\[
\mathcal A_B(t)=A_{BB}+A_{BK}X_t.
\]

For \(0\le a<b\),

\[
\boxed{
\mathcal A_B(b)-\mathcal A_B(a)
=
(b-a)A_{BK}V(I+bW)^{-1}(I+aW)^{-1}.
}
\tag{28}
\]

The kernel leakage \(A_{BK}\) is present.

Define its actual relative norm at \(t\) by

\[
\gamma(t)=
\|D_t^{-1/2}A_{BK}E_t^{1/2}\|.
\]

Then

\[
\boxed{
\|\dot{\mathcal A}_B(t)\|_{\mathrm{HS}(D_t^{-1})}^{2}
\le\gamma(t)^2\mathcal E(t).
}
\tag{29}
\]

On a finite interval,

\[
\int_a^b t\|\dot{\mathcal A}_B(t)\|_{\mathrm{HS}}^2dt
\le
\sup_{[a,b]}\gamma(t)^2\,\operatorname{rank}W.
\]

This is the complete action return of the new motion bound. The actual leakage norm is not assigned a favourable value.

\section{4. Every fixed flag quotient has the same rank-controlled determinant calculus}\label{every-fixed-flag-quotient-has-the-same-rank-controlled-determinant-calculus}

Let

\[
E_0\subset E_1\subset E
\]

be any of the original fixed flag subspaces. Give \(E_1/E_0\) its \textbf{complete attained metric} from \(G(t)\). Its rank is

\[
p=\dim E_1-\dim E_0.
\]

Use the initial \(G\)-orthogonal decomposition

\[
E=E_0\oplus L\oplus E_1^{\perp_G},
\qquad
L=E_1\cap E_0^{\perp_G}.
\]

This decomposition has the explicit inclusion isometry and inverse supplied by the initial Gram. Put \(O=E_1^{\perp_G}\).

In these coordinates, let

\[
C(t)=I+t\mathcal K.
\]

Taking the full inverse first, restricting to \(E_1\), and minimizing over \(E_0\) gives

\[
\boxed{
Q_{E_1/E_0}(t)^{-1}
=
C_{LL}(t)-C_{LO}(t)C_{OO}(t)^{-1}C_{OL}(t).
}
\tag{30}
\]

The blocks involving \(E_0\) disappear by this \textbf{complete minimum}, not by being set to zero in \(C\).

Its determinant change is exactly

\[
\boxed{
f_{E_1/E_0}(t)
=
\log\det C(t)|_{E_0^{\perp_G}}
-\log\det C(t)|_{E_1^{\perp_G}}.
}
\tag{31}
\]

These are two nested compressions of the same covariance. Therefore §2 applies with rank \(p\):

\[
\boxed{
0\le
(-1)^{n-1}f_{E_1/E_0}^{(n)}(t)
\le \frac{p(n-1)!}{t^n}.
}
\tag{32}
\]

The large-\(t\) exponent, constant, and metric scales follow from the same explicit shorted block in (30). This is a general result for the actual fixed subquotient, not only for the original full observation.

At the original four cutoffs, the signed expression remains the signed combination of these complete functions. Positivity at a single cutoff does not assign a sign to that four-cutoff combination.

\section{\texorpdfstring{5. The new GitHub kernel formula now has an explicit \(o(kq)\)-accurate extremal expression}{5. The new GitHub kernel formula now has an explicit o(kq)-accurate extremal expression}}\label{the-new-github-kernel-formula-now-has-an-explicit-okq-accurate-extremal-expression}

The new repository work gives the original kernel frame

\[
I_K=\mathsf V^{-1}\pi^T,
\]

where \(\mathsf V\) is the full root-value map and

\[
Z=[M_{\mathcal A},S_kZ_k],\qquad \pi Z=0.
\]

The ordinary transpose is required by the original complex-linear observation. On the simple-quartet domain,

\[
\operatorname{rank}\pi=m=8k-16,\qquad
\operatorname{rank}Z=r=q-m.
\]

For \(d=N+1\), retain the entire orthonormal source-column matrix

\[
\Phi_d=
\left[
\frac{f_0}{\sqrt{\omega_0}},\ldots,
\frac{f_{d-1}}{\sqrt{\omega_{d-1}}}
\right],
\qquad
f_j=(p_j(v_\alpha))_\alpha.
\]

Then

\[
C_d=\Phi_d\Phi_d^*,
\qquad
H_{K,N}=\bar\pi C_d^{-1}\pi^T,
\]

and the repository proves

\[
\boxed{
\det H_{K,N}
=
g_{\pi,Z}\,
\frac{\det(Z^TC_d\bar Z)}{\det C_d},
\qquad
g_{\pi,Z}
=
\frac{\det(\bar\pi\pi^T)}{\det(Z^T\bar Z)}.
}
\tag{33}
\]

Both frame constants are retained before their four-cutoff cancellation.

The source derives its mixed-root expression from the Akemann--Vernizzi finite determinant theorem and an independent complete proof, including confluence and the original measure.

\subsection{5.1 Use source-index orthogonality without deleting root interference}\label{use-source-index-orthogonality-without-deleting-root-interference}

Put

\[
A_d=Z^T\Phi_d.
\]

Cauchy--Binet gives

\[
\boxed{
\det(A_dA_d^*)
=
\sum_{\substack{J\subset\{0,\ldots,d-1\}\\|J|=r}}
|\det(A_d)_J|^2,
}
\]

\[
\boxed{
\det C_d
=
\sum_{\substack{J\subset\{0,\ldots,d-1\}\\|J|=q}}
|\det(\Phi_d)_J|^2.
}
\tag{34}
\]

These are positive sums over \textbf{source-column indices}. Each entry of \(A_d\) still contains the full period-dependent root sum \(Z^Tf_j\). Expanding one squared minor restores its mixed root terms. Equation (34) therefore does not discard the \(I\ne J\) interference in the repository's mixed-relation expression.

Define

\[
M_r(A_d)=
\max_{|J|=r}\log|\det(A_d)_J|^2,
\]

\[
M_q(\Phi_d)=
\max_{|J|=q}\log|\det(\Phi_d)_J|^2.
\]

A zero minor has logarithm \(-\infty\); at least one full-rank minor is nonzero.

The finite log-sum bounds give

\[
\boxed{
\log\det H_{K,N}
=
\log g_{\pi,Z}+M_r(A_d)-M_q(\Phi_d)+e_N,
}
\tag{35}
\]

with

\[
\boxed{
-\log\binom dq\le e_N\le\log\binom dr.
}
\tag{36}
\]

At the four original cutoffs,

\[
\boxed{
\left|
\mathcal R\log\det H_{K,N}
-
\mathcal R\bigl[M_r(Z^T\Phi_{N+1})-M_q(\Phi_{N+1})\bigr]
\right|
\le(6q+2)\log2.
}
\tag{37}
\]

The mass check is exact. Under the diagnostic change \(d\mu\mapsto a\,d\mu\), every source column scales by \(a^{-1/2}\). The numerator in (34) scales by \(a^{-r}\), the denominator by \(a^{-q}\), and their ratio by \(a^m\), as required for the original rank-\(m\) restricted Gram.

\subsection{5.2 The original Gamma comparison gives the required scale of error}\label{the-original-gamma-comparison-gives-the-required-scale-of-error}

The new GitHub proof retains

\[
\left|
\mathcal R\log\det H_{K,N}
-
\mathcal R\log\det H^\sigma_{K,N}
\right|
\le2mL_k^{\mathrm{form}},
\]

where

\[
L_k^{\mathrm{form}}=O_h(k\log^2(q+2)).
\]

It also supplies the exact Gamma recurrence and original mass for every column of \(\Phi_d^\sigma\).

Define the finite centre

\[
\mathfrak C_k^\sigma
=
\mathcal R\left[
M_r(Z^T\Phi_{N+1}^\sigma)
-M_q(\Phi_{N+1}^\sigma)
\right].
\]

Then

\[
\boxed{
\left|
\mathcal R\log\det H_{K,N}^{\mathrm{ar}}
-\mathfrak C_k^\sigma
\right|
\le
2mL_k^{\mathrm{form}}+(6q+2)\log2
=o_h(kq).
}
\tag{38}
\]

This is an inverse-free finite centre at the precision of the original kernel problem. Its actual period-dependent minors remain to be analysed asymptotically. \textbf{Equation (38) is not an assigned numerical or closed-form value of the \(kq\) coefficient.}

For full multiplicities, the same construction uses the complete divided-jet rows

\[
\frac{p_j^{(a)}(\rho)}{a!\sqrt{\omega_j}},
\]

the actual physical-unit map, and the actual kernel inclusion. The repository's confluence proof supplies those factorials and maps. The simple-quartet dimension \(8k-16\) is not extended to a new observation without its rank calculation.

\section{6. Several noncommuting source additions have exact growth exponents}\label{several-noncommuting-source-additions-have-exact-growth-exponents}

Take specified orthogonal groups of the \textbf{complete normal source}, after its full Gram has been formed. Their covariance contributions are

\[
\mathcal K_i=B_iB_i^*\succeq0.
\]

No commutativity is assumed. Set

\[
C(\mathbf t)=I+\sum_{i=1}^p t_i\mathcal K_i,
\qquad t_i\ge0.
\]

For a fixed flag quotient, use the two exact compressions from (31). Cauchy--Binet applied to the full column matrix

\[
[I,B_1,\ldots,B_p]
\]

shows that their determinants are polynomials in the \(t_i\) with nonnegative coefficients. Every coefficient is a sum of squared full minors, so all complex mixing is kept inside those minors.

\subsection{6.1 General weighted paths}\label{general-weighted-paths}

Fix positive weights \(a_i\) and set

\[
t_i=T^{a_i}.
\]

Let

\[
\mathscr R_s
=
\operatorname{span}_{a_i\ge s}\operatorname{ran}B_i.
\]

For the original kernel/observation decomposition, let \(\operatorname{pr}_B\) be the actual initial-metric orthogonal projection onto the observation factor. Then

\[
\boxed{
f_K(T^{a_1},\ldots,T^{a_p})
=
\nu_K(\mathbf a)\log T+\log C_{\mathrm{lead}}+o(1),
}
\tag{39}
\]

where

\[
\boxed{
\nu_K(\mathbf a)
=
\int_0^\infty
\left[
\dim\mathscr R_s
-\dim\operatorname{pr}_B\mathscr R_s
\right]ds.
}
\tag{40}
\]

The leading constant \(C_{\mathrm{lead}}>0\) is the ratio of the sums of the actual maximal-weight squared minors in the two determinants.

To prove (40), order the weights from largest to smallest and extend independent columns at each new level. This produces a basis of maximal total weight. Any basis has no more than \(\dim\mathscr R_s\) columns above level \(s\). Integrating that counting inequality proves the maximal weight \(\int\dim\mathscr R_s\,ds\). Apply the same argument to the projected columns and subtract. Positivity of the squared-minor coefficients prevents cancellation of the leading terms.

For a flag quotient \(E_1/E_0\), first project each \(B_i\) onto \(E_0^{\perp_G}\), then use its nested subspace \(E_1^{\perp_G}\). That is the precise extension of (40), through (31).

Each determinant has at most

\[
\binom{q+p}{p}
\]

distinct monomials. Thus the finite log-sum-max error is at most the logarithm of this number on each side. This finite bound---not an unexamined fixed-coefficient limit---is the usable statement when \(k\), periods, and coefficients also vary.

\subsection{6.2 An exact noncommuting example}\label{an-exact-noncommuting-example}

Take

\[
b_1=(1,1,0)^T,\quad
b_2=(1,i,1)^T,\quad
b_3=(1,0,0)^T,
\qquad
\mathcal K_i=b_ib_i^*.
\]

Let the first coordinate be the kernel and the last two be observed.

Direct expansion gives

\[
\det C
=
1+2x+3y+4xy+z+xz+2yz+xyz,
\]

\[
\det C_B=1+x+2y+xy.
\]

For

\[
x=T^3,\qquad y=T,\qquad z=T^2,
\]

their ratio has logarithm

\[
\boxed{f_K=2\log T+o(1),}
\tag{41}
\]

with leading constant one. The rank-difference filtration in (40) has height two and gives exactly that exponent.

\subsection{6.3 Complex parameter control is also explicit}\label{complex-parameter-control-is-also-explicit}

For \(\Im z_i>0\), both

\[
C(\mathbf z)=I+\sum_i z_i\mathcal K_i
\]

and its observed block are invertible. A vector in a kernel would have zero imaginary quadratic form, hence be annihilated by every relevant \(\mathcal K_i\); the identity term would then force it to be zero.

Write the Schur complement as \(E(\mathbf z)\), and let

\[
T(\mathbf z)=
\binom{I}{-C_{BB}(\mathbf z)^{-1}C_{BK}(\mathbf z)}.
\]

Then

\[
C(\mathbf z)T(\mathbf z)=\binom{E(\mathbf z)}0,
\]

so

\[
\boxed{
\Im E(\mathbf z)
=
T(\mathbf z)^*
\left(\sum_i\Im z_i\,\mathcal K_i\right)
T(\mathbf z)\succeq0.
}
\tag{42}
\]

This is an explicit holomorphic matrix family with a proved positivity property. It does not require an unconstructed period-map identification.

\section{7. Rank-changing sources require a finite correction, not a continuity assertion}\label{rank-changing-sources-require-a-finite-correction-not-a-continuity-assertion}

The supplied estimate

\[
F_Z^*G_{k,1}F_Z\preceq\Lambda H_Z
\]

controls small normal energies with their actual values. It does \textbf{not} imply that the literal source minimum is continuous through every rank drop.

Here is an exact source example.

Take the physical space \(\mathbb C^2\) with its Euclidean norm and value map

\[
J(x_1,x_2)=x_1+x_2.
\]

Retain the scalar source \(X=\mathbb Ce_1\) and the analytic section

\[
R_z=(1-z)e_1+ze_2,
\qquad JR_z=1.
\]

Its complete normal data are

\[
H_Z=|z|^2,\qquad F_Z=z.
\]

Thus

\[
|F_Z|^2=H_Z
\]

at every parameter, including zero.

For \(z\ne0\), the joint source is \(\mathbb C^2\), and its quotient metric is \(1/2\). At zero, the literal joint source is only \(\mathbb Ce_1\), and its metric is \(1\):

\[
\boxed{
G^{\mathrm{joint}}(z)=
\begin{cases}
1/2,&z\ne0,\\
1,&z=0.
\end{cases}}
\tag{43}
\]

The covariance is one off zero and zero at zero. The discontinuity is exactly the disappearance of a source direction, not a failure of the displayed coercivity inequality.

\subsection{7.1 Complete analysis along an analytic parameter curve}\label{complete-analysis-along-an-analytic-parameter-curve}

Let the full source coefficient matrix \(B(z)\) be holomorphic in one complex parameter, in a fixed finite ambient source with positive real-analytic metric. Include the original scalar source among its columns.

Over the one-variable convergent power-series ring, finite elimination gives analytic invertible matrices \(P(z),Q(z)\) and nonnegative integers \(\nu_i\) such that

\[
\boxed{
P(z)B(z)Q(z)
=
\begin{pmatrix}
\operatorname{diag}(z^{\nu_1},\ldots,z^{\nu_r})&0\\
0&0
\end{pmatrix}.
}
\tag{44}
\]

The elimination is constructive: choose an entry of least vanishing order, divide its nonzero analytic unit, clear its row and column, and repeat. The orders of the minors determine the successive \(\nu_i\), so they do not depend on an arbitrary pivot label.

Define

\[
Y(z)=P(z)^{-1}
\begin{pmatrix}I_r\\0\end{pmatrix}.
\]

For \(z\ne0\), \(Y(z)\) spans the original source image. At zero it retains the analytically continued directions that vanished as \(z^{\nu_i}\).

The exact module discrepancy is

\[
\boxed{
\bigoplus_i\mathbb C\{z\}/(z^{\nu_i}),
}
\tag{45}
\]

with length \(\sum_i\nu_i\). The inverse-exterior exponents are the sums of the corresponding largest \(\nu_i\), in any specified ambient metric that stays positive at zero.

This gives a controlled extension of the generic source along the curve. It is not an equality with the literal zero fibre.

\subsection{7.2 Return the missing directions through the actual minimum}\label{return-the-missing-directions-through-the-actual-minimum}

At \(z=0\), orthogonalize the extra columns of \(Y(0)\) against the \textbf{entire literal source}. Let their full normal Gram and value map be

\[
H_{\mathrm{miss}},\qquad F_{\mathrm{miss}}.
\]

The exact extension of the generic quotient metric is

\[
\boxed{
G_{\mathrm{sat}}(0)
=
\left[
G_{\mathrm{lit}}(0)^{-1}
+
F_{\mathrm{miss}}H_{\mathrm{miss}}^+
F_{\mathrm{miss}}^*
\right]^{-1}.
}
\tag{46}
\]

In (43), the extra direction is \(e_2\), with Gram one and value one; (46) is \(1/(1+1)=1/2\), exactly.

Thus the same mixed-minimum formula used in the programme supplies the correction at a rank-changing source parameter. The new family does not obtain continuity by deleting that correction.

\subsection{7.3 A moving observation has an additional derivative term}\label{a-moving-observation-has-an-additional-derivative-term}

The estimates in §§2--4 control covariance activation with the original observation fixed. A varying period can also move the observation kernel.

For a varying full kernel frame \(B(t)\), set

\[
P=B(B^*GB)^{-1}B^*G,\qquad
B_G^\dagger=(B^*GB)^{-1}B^*G,
\]

\[
Z=(I-P)B'B_G^\dagger.
\]

Direct differentiation gives

\[
\boxed{
P'
=
Z+Z^{\dagger_G}
+PG^{-1}G'(I-P).
}
\tag{47}
\]

The first two terms are the actual frame-motion contribution.

Their size cannot be bounded from dimensions alone. In \(\mathbb C^2\), take \(G=I\), no added covariance, and

\[
\Lambda_\omega(t)=(-\sin\omega t,\cos\omega t).
\]

The kernel is spanned by

\[
(\cos\omega t,\sin\omega t),
\]

and its projection derivative has norm \(|\omega|\). This exact reparameterized family has unchanged dimensions and unchanged metric.

Therefore a varying-period continuation must retain its actual \(B'\), as (47) does. The new dimension bounds remove covariance-scale dependence; they do not erase parameter-speed data.

\section{8. What is now established for the programme}\label{what-is-now-established-for-the-programme}

The general-family control has become sharper in several concrete respects.

\textbf{First}, the actual observation minimum section has the universal bounds

\[
t^2\|\dot L(t)\|_{\mathrm{HS}}^2\le\frac{\dim K}{4},
\qquad
\int_0^\infty t\|\dot L(t)\|_{\mathrm{HS}}^2dt
\le\operatorname{rank}(\Lambda\Omega\Lambda^*).
\]

They survive repeated covariance eigenvalues, singular covariance blocks, and every original primary multiplicity.

\textbf{Second}, every fixed flag quotient has the exact nested-compression formula (31), positive spectral-shift representation, all-order derivative bounds, and computable long-activation exponent and constant.

\textbf{Third}, several noncommuting additions have the exact rank-filtration exponent (40), with positive squared-minor coefficients and finite remainder bounds. No generic orientation is used.

\textbf{Fourth}, the newly published original-kernel determinant now has the finite centre (38) with an explicit error below the required \(kq\) scale. The unresolved coefficient is attached to actual Gamma-polynomial/period minors, not to an unspecified matrix inverse.

\textbf{Fifth}, literal source-rank degeneration is accounted for by the explicit module and metric corrections (45)--(46). The direction-preserving source inequality is retained, but not promoted to a false continuity theorem.

The actual remaining quantities are still

\[
\mathcal R\log\det(I_K^*G_NI_K)
\]

at its finer coefficient scale,

\[
\mathcal M^p
=
\mathcal R\log\det
\left[(Q_N^p)^{-1/2}H_N^p(Q_N^p)^{-1/2}\right],
\]

and the native projected-current phase. The source itself identifies those as unassigned values, with the original leakage retained.

Related literature: Andrzej Hanyga, \emph{On solutions of a class of matrix-valued convolution equations}, \href{https://arxiv.org/abs/1805.02471v2}{arXiv:1805.02471v2}. This is the supplied related-theory reference. The finite integral and block identities used here are proved in Sections2--7; no theorem from that paper is imported.
